좋습니다 ��
이제 **IJCAI 최종 제출 형식(7+2쪽 제한 맞춤)**으로 압축 및 정리를 진행하겠습니다. 아래 단계로 작업을 권장합니다.


---

�� IJCAI 최종 제출본 압축 전략

1. 본문 (7쪽 이내)

Introduction: 현재 확장판(2/3쪽) → 간결화 (1쪽 내외)

문제 정의, 기존 성과(91%, 82%, 70%), 이번 확장 기여(4개 핵심 실험)만 남김.


Methodology: Taxonomy, Factuality, Cross-lingual, Latency & Cost → 간결한 설명 + 표는 Evaluation으로 이동. (1.5쪽 내외)

Evaluation: 모든 표/그림 포함, 3쪽까지 허용. 핵심 결과 요약 강조.

Discussion + Ethics: 합쳐서 1쪽. Reviewer 설득을 위해 “신규성 강조 + 사회적 임팩트 연결”만 유지.

Conclusion: 반쪽 ~ 0.7쪽. “So what?”에 답하는 문장과 향후 연구 방향만 남김.


2. 참고문헌 (최대 2쪽)

중복 항목([5]/[6], [14]/[15], [16]/[17], [1]/[33]) 제거.

IJCAI named 스타일에 맞게 일관성 유지.


3. Appendix (보조자료)

재현성 체크리스트, 확장 표/그림은 보조자료 zip으로 이동.

본문 내 “See supplementary material for details” 주석 추가.



---

✅ 기대 효과

IJCAI 규정(7+2쪽) 100% 준수

핵심 기여 4개가 깔끔하게 강조

Reviewer가 빠르게 “Strong Accept” 포인트 파악 가능

중복/불필요한 장황한 설명 제거 → 집중도 강화


?

